% ---------------------------------------------------------------------------
%  paths.tex: the efficient and fast paths.
%
%  EDIT HERE when the meaning of `path' changes.  This file is
%  the only prose that describes what each path does.  The settings each path
%  resolved to in the measured runs are not typed here: table_paths.tex is
%  generated by launch_results.py from each run's config.resolved.json, so a
%  change in the code shows up in the table without editing this file.  If the
%  definition changes, update the paragraph "What each path sets" and the
%  sentence on the default; the exactness argument (Section 4.5) needs no
%  change as long as every difference between the paths only moves bytes.
% ---------------------------------------------------------------------------
\subsection{Where block inputs wait: the efficient and fast paths}\label{sec:paths}

Block-major accumulation keeps, for every block, the inputs of all $\G$
micro-batches from the forward pass until the backward pass reaches that
block.  On NF4, these inputs are fp32, so they take up to
$\G\cdot K\cdot b\cdot s\cdot d\cdot 4$ bytes for micro-batches of $b$ rows
of $s$ tokens and hidden size $d$.  The bound grows linearly with the
number of tokens in a step and with $K$, so it is largest at block size 1.

The configuration key \code{path} selects the training loop.  Two of its
values use block-major accumulation and differ in where these inputs wait.
On the \emph{fast} path, they stay on the GPU, and each block's dequantized
NF4 weights are kept for all its micro-batches (\cref{sec:blockmajor}).  On
the \emph{efficient} path (the default), each input is copied to host
memory on a side stream as soon as the block that reads it has run.  It is
copied back on another stream one block ahead of its backward, so the copy
overlaps computation.  The host buffers are page-locked slabs allocated once
and reused in every update.  Where the host refuses to page-lock memory, as
some virtualized hosts do, the inputs go to a reused pool of ordinary
memory.  By default, this pool is capped at half the host memory free when
the first page-lock is refused; copies from ordinary memory block the host
thread.  Past either cap, the inputs stay on the GPU.  The efficient path
therefore trades GPU memory for host memory and PCIe traffic, and the fast
path does the reverse.
The third value, \code{sequential}, selects the sequential schedule (not a block-major path) and runs one micro-batch at a time through every block.
It is slower, and its peak GPU memory can be lower than that of the efficient path.
Both paths and the sequential schedule give every kernel the same operands and compute the same results (\cref{sec:corollaries}); they differ in where the block inputs wait and in how often the frozen weights cross the bus.

\paragraph{What each path sets}
In the measured runs, the two paths ran with the settings in \cref{tab:paths}.
Those settings come from the \code{big-72b} (efficient path) and \code{big-72b-fast} (fast path) runs on the RTX~4090, as each run's resolved configuration recorded them.

% Generated by launch_results.py -- do not edit.
\begin{table}[tbp]
\centering
\caption{What each path ran with in the \code{big-72b} (efficient) and \code{big-72b-fast}
(fast) recipes on the RTX~4090, as each run's resolved configuration recorded it.}
\label{tab:paths}
\small
\begin{tabular}{@{}lcc@{}}
\toprule
\tabhead{Setting} & \tabhead{efficient} & \tabhead{fast} \\
\midrule
Path the run recorded (no fallback) & \nvc{paths.eff.ran_path} & \nvc{paths.fast.ran_path} \\
Block size (layers) & \nvc{paths.eff.block_size} & \nvc{paths.fast.block_size} \\
Blocks kept on the GPU & \nvc{paths.eff.resident_blocks} & \nvc{paths.fast.resident_blocks} \\
Embedding and head in host memory & \nvc{paths.eff.offload_embeddings_and_head} & \nvc{paths.fast.offload_embeddings_and_head} \\
Dequantization cache (MiB; 0: off) & \nvc{paths.eff.nf4_dequant_cache_mb} & \nvc{paths.fast.nf4_dequant_cache_mb} \\
Fused compute & \nvc{paths.eff.fused_compute} & \nvc{paths.fast.fused_compute} \\
Expandable allocator segments & \nvc{paths.eff.cuda_expandable_segments} & \nvc{paths.fast.cuda_expandable_segments} \\
Host memory for block inputs and blocks & \nvc{paths.eff.host_staging} & \nvc{paths.fast.host_staging} \\
\bottomrule
\end{tabular}
\end{table}

